% Article VII: pre-existing proofs incorporated into the body.
% Complete source: fuentes_conservadas/17_rebote_einstein_cartan.tex.
% Includes the scaling law and the proofs of bounce and persistence.
% The amplitude s_0 is an output of the preceding material realization;
% this fragment alone does not certify evaluation of that realization.

\section{Memory dilution and persistence of the bounce}
\label{vii:sec:dilucion-rebote}

Memory evolution and spatial expansion are two readouts of
the same sequence of states. The refinement index can be eliminated
between them; this elimination determines the dilution exponent that the
cosmological realization transmits to its torsional contribution.
The initial amplitude depends on evaluation of the specific material
current in its variational response. The scaling law and bounce
theorem below preserve this dependence.

Theorem~\ref{vii:thm:amplitud-material-explicita} determines this
amplitude as a functional of the material state and comoving volume.
Corollary~\ref{vii:cor:dominio-positivo-conservado} provides
an explicit domain of states with positive, conserved amplitude.
In this domain, evaluation of \(s_0\) precedes the cosmological
integration developed below. The ordinary density
\(\rho_0>0\) and its barotropic law are specified additionally in
Section~\ref{vii:sec:umbral-rebote}: the amplitude-positivity witness
does not by itself evaluate this material density.

\subsection{The \texorpdfstring{\(a^{-6}\)}{a to the power -6} dilution law}
\label{vii:sec:dilucion}

\begin{proposition}[Elimination of the memory index]
\label{vii:prop:dilucion}
Let \(a_{\mathrm{ref}}>0\), \(\varphi>1\), and let the route readers be
\[
 \frac{a_n}{a_{\mathrm{ref}}}=\varphi^{n/3},\qquad
 M_n=\varphi^{-2n}.
\]
Then
\[
 M_n=\left(\frac{a_n}{a_{\mathrm{ref}}}\right)^{-6}
\]
for every index in the sequence. If the quadratic-density realization
is \(s_n^2=s_*^2M_n\), with \(s_*^2>0\), then
\[
 s_n^2=s_*^2\left(\frac{a_n}{a_{\mathrm{ref}}}\right)^{-6}.
\]
\end{proposition}

\begin{proof}
Raising the first equality to the power \(-6\) gives
\(\varphi^{-2n}\), which is exactly the second reader. Multiplication
by the initial amplitude yields the second statement. The exponent
is determined by the ratio of the exponents of the two readers;
the amplitude is the evaluation of the reference state.
\end{proof}

The continuous homogeneous realization preserves this constitutive law:
the torsional density is written \(s_0a^{-6}\), where the normalization
of the factor \(a\), the material coefficient and the reference amplitude
have been absorbed into \(s_0\). The scaling law and evaluation of
\(s_0\) are successive operations of the construction.

\subsection{Homogeneous threshold and transversality}
\label{vii:sec:umbral-rebote}

Consider the spatially flat homogeneous chart, in units \(c=1\),
with constant \(G>0\), positive densities \(\rho_0,s_0\), and constant
barotropic parameter \(w<1\). The specified realization satisfies
\(\Lambda_{\mathrm{eff}}=0\) and the equations
\begin{align}
 \rho_{\mathrm m}(a)&=\rho_0a^{-3(1+w)},&
 \rho_{\mathrm s}(a)&=s_0a^{-6},\\
 H^2&=\frac{8\pi G}{3}
        \bigl(\rho_{\mathrm m}-\rho_{\mathrm s}\bigr)
        =:F_0(a),&
 \dot H&=-4\pi G
        \bigl((1+w)\rho_{\mathrm m}-2\rho_{\mathrm s}\bigr).
 \label{vii:eq:dinamica-homogenea}
\end{align}
The couplings belong to the dimensional section already constructed.
The expression for \(\rho_{\mathrm s}\) incorporates the sign of its
contribution in this gravitational realization.

\begin{theorem}[Existence and uniqueness of the transverse threshold]
\label{vii:thm:rebote}
There is a unique \(a_{\mathrm b}>0\) with \(F_0(a_{\mathrm b})=0\):
\[
 a_{\mathrm b}^{\,3(1-w)}=\frac{s_0}{\rho_0}.
\]
The region permitted by \(H^2\geq0\) is \(a\geq a_{\mathrm b}\).
At the threshold, with
\(\rho_{\mathrm b}=\rho_{\mathrm m}(a_{\mathrm b})
                  =\rho_{\mathrm s}(a_{\mathrm b})>0\),
\[
 \dot H_{\mathrm b}=4\pi G(1-w)\rho_{\mathrm b}>0.
\]
The solution crossing \(H=0\) has a strict minimum of \(a\),
with local expansion
\[
 a(t)=a_{\mathrm b}\left[
 1+2\pi G(1-w)\rho_{\mathrm b}(t-t_{\mathrm b})^2
 +O((t-t_{\mathrm b})^3)\right].
\]
\end{theorem}

\begin{proof}
The factorization
\[
 F_0(a)=\frac{8\pi G}{3}a^{-6}
             \bigl(\rho_0a^{3(1-w)}-s_0\bigr)
\]
reduces its zeros to that of a strictly increasing function from
\((0,\infty)\) onto \((-s_0,\infty)\). This proves existence,
uniqueness and the sign of \(F_0\) on either side of the threshold.
The second equation in \eqref{vii:eq:dinamica-homogenea}, evaluated
at equality of the densities, gives the formula for \(\dot H_{\mathrm b}\).
Since \(\dot a=aH\), at the threshold
\(\ddot a_{\mathrm b}=a_{\mathrm b}\dot H_{\mathrm b}>0\).
The equations are smooth for \(a>0\), so their Taylor expansion
gives the stated expression. The change of sign of \(H\) is
transverse: it passes from contraction to expansion.
\end{proof}

For \(w=1\), both terms have the same exponent, and equality of
amplitudes determines a degenerate case. For \(w>1\), the sign of
\(\dot H\) at a coincidence of densities is negative.
The theorem's domain \(w<1\) expresses precisely the order between
the two dilution exponents.

\subsection{Local persistence under controlled perturbations}
\label{vii:sec:persistencia}

Persistence of the threshold uses two independent estimates:
control of the sign at the endpoints and of the derivative on the
interval. Let \(I=[a_-,a_+]\subset(0,\infty)\), with
\(a_-<a_{\mathrm b}<a_+\), be chosen so that
\[
 F_0(a_-)<0<F_0(a_+),\qquad
 m_I:=\inf_{a\in I}F_0'(a)>0.
\]
Such an interval exists because
\[
 F_0'(a_{\mathrm b})
   =\frac{8\pi G}{a_{\mathrm b}}(1-w)\rho_{\mathrm b}>0.
\]
Define
\(\delta_I=\min\{-F_0(a_-),F_0(a_+)\}>0\).

\begingroup\fontsize{10}{11.7}\selectfont
\setlength{\parskip}{2pt}\setlength{\abovedisplayskip}{5pt}\setlength{\belowdisplayskip}{5pt}
\begin{theorem}[Local persistence of the transverse bounce]
\label{vii:thm:persistencia}
Let \(R\in C^1(I)\) be a perturbation satisfying
\[
 \|R\|_{\infty,I}<\delta_I,\qquad
 \|R'\|_{\infty,I}<m_I.
\]
Then \(F_R=F_0+R\) has a unique zero \(a_R\in(a_-,a_+)\), and
\(F_R'(a_R)>0\). The realization \(H^2=F_R(a)\) admits a local
transverse bounce, with
\[
 \dot H_R=\frac{a_R}{2}F_R'(a_R)>0
\]
at the threshold. If \(R(a,\eta)\) is jointly \(C^1\), the persistent
zeros depend locally in a \(C^1\) manner on the parameter \(\eta\).
\end{theorem}

\begin{proof}
The first bound preserves
\(F_R(a_-)<0<F_R(a_+)\); the intermediate value theorem provides
an interior zero. The second gives
\(F_R'(a)\geq m_I-\|R'\|_{\infty,I}>0\) on \(I\), so the zero
is unique and simple.

For \(a>a_R\) close to the zero,
\(F_R(a)=F_R'(a_R)(a-a_R)+o(a-a_R)\).
The integral
\[
 t-t_{\mathrm b}=\int_{a_R}^{a}
          \frac{d\xi}{\xi\sqrt{F_R(\xi)}}
\]
is finite and strictly increasing. Its inverse defines the expanding
branch; time reflection defines the contracting branch.
They join with \(\dot a(t_{\mathrm b})=0\) and
\[
 a(t)-a_R
 =\frac{a_R^2F_R'(a_R)}4(t-t_{\mathrm b})^2
   +o((t-t_{\mathrm b})^2).
\]
Away from the threshold, differentiating \(H^2=F_R(a)\) gives
\(\dot H=aF_R'(a)/2\); its continuous limit gives the positive
expression at \(a_R\). The \(C^1\) regularity of \(F_R\) ensures
continuity of the acceleration at this junction. Finally,
\(\partial_aF_R(a_R)>0\) permits application of the implicit function
theorem when the perturbation depends jointly on \((a,\eta)\).
\end{proof}

The first bound controls existence; the second ensures local
uniqueness and transversality. The stability proved here concerns
this threshold in the selected interval. Torsional amplitude enters
its location, whereas the relation between exponents determines
the temporal orientation of the crossing.

\par\endgroup
